\section{第13组：光的干涉}
\sourceproblem{13-T1}{79}{83}
\problemtext 同频同初相相干光在折射率n的均匀介质内传播，两路几何路程差为 $r_2-r_1$，求光程差和相位差。
\workings $\Delta L=n(r_2-r_1)$，传播相位滞后差 $\Delta\phi=2\pi\nu\Delta L/c$。若定义同一时刻第二束的相位减第一束，则为 $-2\pi\nu\Delta L/c$。
\crosscheck 真空波长 $\lambda_0=c/\nu$；不重复乘折射率。
\sourceproblem{13-T2}{79}{83}
\problemtext 600 nm光由空气斜入射折射率1.23、厚1 cm平板再出射，入射角30度，板外两段各5 cm，求介质中频率、速率、波长及总几何路程与光程。
\workings $\nu=c/\lambda_0=5.00\times10^{14}$ Hz，$u=c/1.23\approx2.439\times10^8$ m/s，$\lambda=600/1.23\approx487.8$ nm。折射角满足 $\sin r=0.5/1.23$，板内路径 $L=1/\sqrt{1-(0.5/1.23)^2}$ cm，约1.0945 cm。几何路程 $10+L\approx11.095$ cm，光程 $10+1.23L\approx11.346$ cm。
\crosscheck 折射前后频率不变，板内路径并非板厚；空气折射率近似为1。
\sourceproblem{13-T3}{79}{83}
\problemtext 双缝间距0.30 mm、屏距1.20 m，第五暗纹距中央11.39 mm，求波长。
\workings 第五暗纹对应 $(5-1/2)\lambda$。近轴 $\lambda=dx/(4.5D)\approx6.328\times10^{-7}$ m，即约633 nm。用远场角度修正 $\lambda=d x/[4.5\sqrt{D^2+x^2}]$，结果约632.75 nm。
\crosscheck 第一暗纹对应半波长而非一波长；近轴角约0.00949 rad。
